{\Large\textbf{Solenoid}}

A solenoid of length $l = 20\ cm$ is wound around a vertical, cylindrical test
tube made of glass, filled with water. The solenoid is thermally insulated from
the water. The height of the water level is approximately $20\ cm$ above the
upper end of the solenoid, the diameter of the test tube is $1\ cm$, the number
of turns of the coil is $N = 6000$.

\begin{center}\includegraphics[width=0.25\linewidth]{figures/EuPhO_2018_Q2_fig1.png}\end{center}

The atmospheric pressure is $p_0 = 101\ kPa$, the temperature of water is
$293\ K$. Magnetic susceptibility of water is
$\chi \equiv \mu_r\ - 1 = -9.04 \cdot 10^{-6}$. Vacuum permeability
$\mu_0 = 12.57 \cdot 10^{-7}\ \mathrm{H/m}$. The current in the solenoid is
slowly increased until the water starts boiling inside the coil. At which
current strength would this happen? Make reasonable approximations when needed.
Note that the required current strength might be slightly too large for the
present technology.
